% !TEX root = ../ApproxDA-TransUNet.tex
\section{Results}
\label{sec:experiments}

We evaluate ApproxDA-TransUNet on five benchmarks spanning tasks with fundamentally different spatial context requirements, comparing against DA-TransUNet and other competitive baselines.

\subsection{Experimental Setup}

\textbf{Datasets.} We evaluate on five representative medical image segmentation benchmarks.

\textit{Synapse Multi-Organ CT}~\cite{landman2015miccai}: 30 abdominal CT scans with an 18/12 train/test split covering 8 organ classes. Metrics: mean DSC and HD95.

\textit{ACDC}~\cite{bernard2018deep}: 100 short-axis cardiac MRI studies with a 70/10/20 train/val/test split covering 3 cardiac structures (RV, Myo, LV). Metrics: per-class and mean DSC, HD95.

\textit{Kvasir-SEG}~\cite{jha2020kvasir}: 1000 endoscopy polyp images (binary) with an 800/200 train/test split. Metrics: DSC and mIoU.

\textit{CVC-ClinicDB}~\cite{bernal2015cvc}: 612 colonoscopy polyp images (binary) with a 489/123 train/test split. Metrics: DSC, mIoU, and HD95.

\textit{ISIC 2018}~\cite{codella2019isic}: 2594 dermoscopy skin lesion images (binary) with a 2075/519 train/test split. Metrics: DSC and mIoU.

\textbf{Implementation Details.} All experiments use an R50-ViT-B/16 backbone pretrained on ImageNet-21K, trained on NVIDIA Tesla T4 (16\,GB) GPUs. Optimizer: SGD, lr$=$0.01; batch size 24; input $224{\times}224$; 300 epochs; best checkpoint selected by evaluating every 15 epochs. Default ApproxDA-TransUNet configuration: $M{=}7$, $r{=}32$, $G{=}8$.

\textbf{Gate Protocol.} ApproxDA-TransUNet supports four gate modes: \texttt{pam} ($g{=}1.0$, PAM only), \texttt{cam} ($g{=}0.0$, CAM only), \texttt{fixed} ($g{=}0.5$, equal blend), and \texttt{learn} ($g$ learned per channel). Each experiment's gate configuration is specified in the corresponding table footnote.

\subsection{Comparison with State-of-the-Art}

Tables~\ref{tab:synapse}--\ref{tab:isic} compare ApproxDA-TransUNet against published baselines across all five benchmarks.

\textbf{Synapse.} Table~\ref{tab:synapse} benchmarks against 11 published methods. CNN-based models plateau at 76--78\% DSC, with Att-UNet leading at 77.77\%. Transformer-based methods progressively improve: TransUNet (77.48\%), UCTransNet (78.23\%), MIM (78.59\%), Swin-Unet (79.13\%), and DA-TransUNet (79.80\%). ApproxDA-TransUNet with gate=pam, $M{=}28$, $r{=}32$ achieves \textbf{80.94\%} DSC and 27.49\,mm HD95---the highest DSC across all compared methods, surpassing DA-TransUNet by $+$1.14\%. The intermediate window $M{=}28$ captures sufficient cross-organ context while limiting spurious long-range correlations; Section~\ref{sec:understanding} analyzes this behavior in detail.

\begin{table}[ht]
  \caption{Segmentation Results on Synapse Multi-Organ CT}
  \label{tab:synapse}
  \begin{center}
  \renewcommand{\arraystretch}{1.15}
  \setlength{\tabcolsep}{6pt}
  \begin{tabular}{lcc}
    \toprule
    Method & DSC (\%) $\uparrow$ & HD95 (mm) $\downarrow$ \\
    \midrule
    \multicolumn{3}{l}{\textit{CNN-based}} \\
    U-Net~\cite{ronneberger2015unet}              & 76.85 & 39.70 \\
    U-Net++~\cite{zhou2018unetpp}                 & 76.91 & 36.93 \\
    Residual U-Net~\cite{diakogiannis2020resunet} & 76.95 & 38.44 \\
    MultiResUNet~\cite{ibtehaz2020multiresunet}   & 77.42 & 36.84 \\
    Att-Unet~\cite{oktay2018attention}            & 77.77 & 36.02 \\
    \midrule
    \multicolumn{3}{l}{\textit{Transformer-based}} \\
    TransUNet~\cite{chen2021transunet}            & 77.48 & 31.69 \\
    UCTransNet~\cite{wang2022uctransnet}          & 78.23 & 26.75 \\
    TransNorm~\cite{azad2022transnorm}            & 78.40 & 30.25 \\
    MIM~\cite{wang2022mim}                        & 78.59 & 26.59 \\
    Swin-Unet~\cite{cao2022swinunet}              & 79.13 & \textbf{21.55} \\
    DA-TransUNet~\cite{sun2024datransunet}        & \underline{79.80} & \underline{23.48} \\
    \midrule
    \multicolumn{3}{l}{\textit{Dual-attention approximation (ours)}} \\
    \textbf{ApproxDA-TransUNet (ours)}     & \textbf{80.94} & 27.49 \\
    \bottomrule
  \end{tabular}
  \end{center}
  \vspace{-0.3em}
  \footnotesize{ApproxDA: gate=pam, $M{=}28$, $r{=}32$, $G{=}8$.}
\end{table}

\textbf{ACDC.} On cardiac MRI segmentation (Table~\ref{tab:acdc}), ApproxDA-TransUNet with gate=pam, $M{=}7$ achieves a mean DSC of 88.97\%---below TransUNet (89.71\%) and Swin-Unet (90.00\%), but with the \textbf{highest} per-class DSC on both RV (\textbf{89.61\%}) and Myo (\textbf{85.91\%}). The class-level pattern is consistent with the inductive-bias interpretation: local windowed attention benefits locally-resolvable structures whose boundaries are determined by fine-grained intensity gradients (crescent-shaped RV; thin-ring Myo), while the LV cavity (91.40\%) degrades 3--4\,pp below all baselines. LV delineation requires estimating global cavity extent---a property that spans multiple M=7 windows and cannot be recovered from within-window context alone. Within a nominally low-GCS dataset ($\Delta\text{DSC}{=}0.73$\,pp overall), the locality prior therefore produces opposing sub-class effects: improving locally-bounded structures while hurting globally-defined ones. We discuss this further in Section~\ref{sec:conclusion}.

\begin{table}[ht]
  \caption{Segmentation Results on ACDC Cardiac MRI}
  \label{tab:acdc}
  \begin{center}
  \renewcommand{\arraystretch}{1.15}
  \setlength{\tabcolsep}{2pt}
  \resizebox{\columnwidth}{!}{\footnotesize\begin{tabular}{lcccc}
    \toprule
    Method & RV (\%) $\uparrow$ & Myo (\%) $\uparrow$ & LV (\%) $\uparrow$ & Mean (\%) $\uparrow$ \\
    \midrule
    U-Net~\cite{ronneberger2015unet}$^\dag$          & 87.10 & 80.63 & 94.92 & 87.55 \\
    Att-UNet~\cite{oktay2018attention}$^\dag$         & 87.58 & 79.20 & 93.47 & 86.75 \\
    TransUNet~\cite{chen2021transunet}                & \underline{88.86} & 84.53 & \underline{95.73} & \underline{89.71} \\
    Swin-Unet~\cite{cao2022swinunet}                  & 88.55 & \underline{85.62} & \textbf{95.83} & \textbf{90.00} \\
    DA-TransUNet~\cite{sun2024datransunet}$^\ddag$   & \multicolumn{4}{c}{\textit{pending re-run}} \\
    \midrule
    \shortstack[l]{\textbf{ApproxDA-TransUNet}\\\textbf{(ours)}} & \textbf{89.61} & \textbf{85.91} & 91.40 & 88.97 \\
    \bottomrule
  \end{tabular}}
  \end{center}
  \vspace{-0.3em}
  {\footnotesize $^\dag$R50-backbone variant as reported in~\cite{chen2021transunet}. $^\ddag$DA-TransUNet not evaluated on ACDC in original paper; pending re-run under our protocol. ApproxDA: gate=pam, $M{=}7$, $r{=}32$, $G{=}8$.}
\end{table}

\textbf{Kvasir-SEG.} On polyp segmentation (Table~\ref{tab:kvasir}), U-Net leads the CNN-based group (88.22\% DSC) while TransUNet (87.91\%) and DA-TransUNet$^\dag$ (88.44\%) are the strongest Transformer-based entries. ApproxDA-TransUNet with gate=pam, $M{=}56$ achieves \textbf{90.17\%} DSC, \textbf{84.27\%} mIoU, and \textbf{44.35\,mm} HD95---state-of-the-art across all three metrics, with a $+$1.73\% DSC gain and $-$8.69\,mm HD95 improvement over the nearest baseline.

\begin{table}[ht]
  \caption{Segmentation Results on Kvasir-SEG Polyp}
  \label{tab:kvasir}
  \begin{center}
  \setlength{\tabcolsep}{3pt}
  \begin{tabular}{lccc}
    \toprule
    Method & DSC (\%) $\uparrow$ & mIoU (\%) $\uparrow$ & HD95 (mm) $\downarrow$ \\
    \midrule
    U-Net~\cite{ronneberger2015unet}           & 88.22 & 80.12 & -- \\
    Att-Unet~\cite{oktay2018attention}         & 86.61 & 78.01 & -- \\
    U-Net++~\cite{zhou2018unetpp}              & 86.57 & 77.67 & -- \\
    Res-UNet~\cite{diakogiannis2020resunet}    & 77.85 & 66.04 & -- \\
    TransUNet~\cite{chen2021transunet}         & 87.91 & 80.03 & -- \\
    DA-TransUNet~\cite{sun2024datransunet}$^\dag$ & \underline{88.44} & \underline{81.70} & \underline{53.04} \\
    \midrule
    \textbf{ApproxDA-TransUNet (ours)}         & \textbf{90.17} & \textbf{84.27} & \textbf{44.35} \\
    \bottomrule
  \end{tabular}
  \end{center}
  \vspace{-0.3em}
  \footnotesize{$^\dag$We are able to closely reproduce paper-reported results (88.47\% DSC, 81.02\% mIoU) and capture the HD95. ApproxDA: gate=pam, $M{=}56$, $r{=}32$.}
\end{table}

\textbf{CVC-ClinicDB.} On colonoscopy polyp segmentation (Table~\ref{tab:cvc}), DA-TransUNet$^\ddag$ achieves 89.47\% DSC and 82.51\% mIoU---the strongest published baseline on this dataset. ApproxDA-TransUNet with gate=pam, $M{=}112$ achieves \textbf{90.99\%} DSC, \textbf{85.09\%} mIoU, and \textbf{14.77\,mm} HD95---surpassing the nearest baseline by $+$1.52\% DSC and $+$2.58\% mIoU. The DSC span across $M{\in}\{7,28,56,112\}$ is 0.62\,pp, confirming CVC-ClinicDB is window-robust---consistent with its binary task structure and low GCS, as analyzed in Section~\ref{sec:understanding}.

\begin{table}[ht]
  \caption{Segmentation Results on CVC-ClinicDB Polyp}
  \label{tab:cvc}
  \begin{center}
  \setlength{\tabcolsep}{3pt}
  \begin{tabular}{lccc}
    \toprule
    Method & DSC (\%) $\uparrow$ & mIoU (\%) $\uparrow$ & HD95 (mm) $\downarrow$ \\
    \midrule
    U-Net~\cite{ronneberger2015unet}                 & 86.93 & 78.21 & -- \\
    Att-Unet~\cite{oktay2018attention}               & 87.41 & 79.35 & -- \\
    U-Net++~\cite{zhou2018unetpp}                    & 87.14 & 78.47 & -- \\
    Res-UNet~\cite{diakogiannis2020resunet}          & 74.22 & 59.02 & -- \\
    TransUNet~\cite{chen2021transunet}               & 89.01 & 81.63 & -- \\
    DA-TransUNet~\cite{sun2024datransunet}           & \underline{89.47} & \underline{82.51} & -- \\
    \midrule
    \textbf{ApproxDA-TransUNet (ours)}               & \textbf{90.99} & \textbf{85.09} & \textbf{14.77} \\
    \bottomrule
  \end{tabular}
  \end{center}
  \vspace{-0.3em}
  \footnotesize{ApproxDA: gate=pam, $M{=}112$, $r{=}32$.}
\end{table}

\textbf{ISIC 2018.} On skin lesion segmentation (Table~\ref{tab:isic}), CNN-based methods cluster between 83--88\% DSC, with TransUNet (88.78\%) and DA-TransUNet (88.88\%) leading among Transformer-based approaches. ApproxDA-TransUNet with gate=learn, $M{=}7$ achieves \textbf{89.58\%} DSC ($+$0.70\% over DA-TransUNet) and 82.66\% mIoU, attaining the best DSC across all compared methods. This window-robust pattern is consistent with the low-GCS behavior observed on Kvasir-SEG (GCS defined in Section~\ref{sec:obs2}).

\begin{table}[ht]
  \caption{Segmentation Results on ISIC 2018 Skin Lesion}
  \label{tab:isic}
  \begin{center}
  \renewcommand{\arraystretch}{1.15}
  \setlength{\tabcolsep}{6pt}
  \begin{tabular}{lcc}
    \toprule
    Method & DSC (\%) $\uparrow$ & mIoU (\%) $\uparrow$ \\
    \midrule
    U-Net~\cite{ronneberger2015unet}           & 87.22 & 81.14 \\
    Att-Unet~\cite{oktay2018attention}         & 87.60 & 81.51 \\
    U-Net++~\cite{zhou2018unetpp}              & 87.30 & 81.33 \\
    Res-UNet~\cite{diakogiannis2020resunet}    & 83.32 & 76.51 \\
    TransUNet~\cite{chen2021transunet}         & 88.78 & 82.63 \\
    DA-TransUNet~\cite{sun2024datransunet}     & \underline{88.88} & \textbf{82.78} \\
    \midrule
    \textbf{ApproxDA-TransUNet (ours)}         & \textbf{89.58} & \underline{82.66} \\
    \bottomrule
  \end{tabular}
  \end{center}
  \vspace{-0.3em}
  \footnotesize{ApproxDA: gate=learn, $M{=}7$, $r{=}32$, $G{=}8$.}
\end{table}

Across five benchmarks (Figure~\ref{fig:qualitative}), ApproxDA-TransUNet achieves state-of-the-art or competitive results. The pattern is consistent with GCS theory: large gains on high-GCS Synapse ($+$1.14\%) and low-GCS texture-boundary tasks (Kvasir $+$1.73\%, CVC $+$1.52\%, ISIC $+$0.70\%); comparable performance on low-GCS ACDC where concentric ring structure provides neither cross-structure context benefit nor texture-prior alignment.

\begin{figure}[b]
  \centering
  \includegraphics[width=\columnwidth]{figures/fig_qualitative_cross_task.png}
  \caption{Qualitative segmentation comparison (columns: input, ground truth, DA-TransUNet, ApproxDA). \textbf{Top (Synapse, $M{=}28$):} multi-organ CT; ApproxDA 80.94\% DSC ($+$1.14\%). \textbf{Middle (Kvasir-SEG, $M{=}56$):} polyp; 90.17\% DSC ($+$1.73\%). \textbf{Bottom (ISIC 2018, $M{=}7$):} skin lesion; 89.58\% DSC ($+$0.70\%).}
  \label{fig:qualitative}
\end{figure}

\subsection{Efficiency Analysis}

Although the proposed approximation substantially reduces the theoretical complexity of the dual-attention module, the overall computational cost is dominated by the ViT backbone and convolutional encoder, which remain unchanged. As shown in Table~\ref{tab:efficiency}, end-to-end GFLOPs and memory exhibit only marginal differences from DA-TransUNet. The practical advantage of ApproxDA-TransUNet therefore lies in improved architectural compatibility rather than total model complexity reduction: DA-TransUNet does not support DDP training, while ApproxDA-TransUNet trains natively across multiple GPUs, enabling cost-effective distributed training on commodity clusters.

\begin{table*}[ht]
  \caption{Efficiency Comparison on Synapse (NVIDIA Tesla T4)}
  \label{tab:efficiency}
  \begin{center}
  \setlength{\tabcolsep}{5pt}
  \begin{tabular}{lcccccc}
    \toprule
    Method & Params (M) & GFLOPs & Train VRAM & Train Time & Infer (s/vol) & Multi-GPU \\
    \midrule
    DA-TransUNet~\cite{sun2024datransunet} & 107.95 & 30.2 & 11.5\,G & 11.4h & 120.0 & No$^\dagger$ \\
    \midrule
    ApproxDA-TransUNet (gate=pam, $M$=7)  & 112.98 & 32.1 & 6.4\,G ($\times$2 GPU) & 8.3h ($\times$2 GPU) & 121.3 & Yes (DDP) \\
    ApproxDA-TransUNet (gate=learn, $M$=7) & 114.90 & 32.1 & 10.6\,G & 11.5h & 114.2 & Yes (DDP) \\
    \bottomrule
  \end{tabular}
  \end{center}
  \vspace{-0.3em}
  \footnotesize{$^\dagger$DA-TransUNet does not support multi-GPU training. ApproxDA-TransUNet uses Distributed Data Parallel (DDP); per-GPU batch 12, LR unchanged.}
\end{table*}
